%% P03 --- Orders of magnitude: O-notation, FLOPs and memory
%%
%% Stub, generated by tools/gen_stubs.py from tools/programs.json.
%% The brief below is the contract for this program: what it must argue, and
%% what it must buy the reader. Writing the program means deleting the
%% \programstub{} block and replacing it with the program -- after which this
%% file is never regenerated. If the brief turns out to be wrong once the
%% mathematics has been worked, change it in the manifest and record the
%% contradiction in CLAUDE.md under *Resolved questions*.

\program{Orders of magnitude: O-notation, FLOPs and memory}
\label{prog:P03}

\programstub{%
Argument: asymptotic notation is a statement about growth and says nothing
about which of two implementations is faster today; both facts matter.
Payoff: bytes, FLOPs as a rate, and arithmetic intensity as a ratio, which
is the whole justification for kernel fusion; why a large matrix multiply is
compute-bound and an elementwise operation is memory-bound. DEPENDENCY NOTE:
an earlier brief had the reader count the parameters of a transformer block
here, three parts before any program has said what a matrix is. That count
moves to P32, where the shapes exist. What stays is magnitude and cost,
which need only arithmetic and sigma notation.

\medskip
\textbf{Planned length:} 45 frames.
}
